\chapter*{Etterord: eit fag som veks}
\addcontentsline{toc}{chapter}{Etterord: eit fag som veks}
\markboth{Etterord}{}

\begin{center}\itshape
Ikkje ferdig er ikkje det same som utan grunn.
\end{center}
\vspace{8mm}

Eg har ført forsvaret så sterkt eg kunne, og no, på slutten, skuldar eg lesaren det ærlege ordet. Designfaget er ikkje ferdig. Det har ikkje den ferdige skikkelsen til dei eldste profesjonane; det har ikkje éin samla etikk som alle utøvarar svarar til, og det har inga handheving med tenner. Dette er den hardaste delen av saka i \textsc{Grunnlaust}, og eg vil ikkje gå rundt henne. Eg vil møte henne ærleg, og så seie kvifor ho likevel ikkje provar det ho skal prove. For skilnaden mellom eit fag utan fundament og eit ungt fag som framleis modnar, er heile skilnaden — og alt eg har skrive i denne boka, kviler på at design er det andre og ikkje det fyrste.

\section{Det sanne i kritikken}

Lat meg fyrst gje motparten det som er motparten sitt. Det sterkaste i \textsc{Grunnlaust} er ikkje påstanden om at design manglar metode, for det viste boka her at det ikkje gjer. Det er ikkje påstanden om at faget manglar ei kunnskapsform, for «designerly ways of knowing» er ein reell og granska tradisjon. Det sterkaste er noko snevrare og hardare: at design manglar \emph{bindande handheving} og \emph{éin samla etikk}. Og det er sant. Ein lege kan miste retten til å praktisere; ein advokat kan bli utestengd frå standen; ein revisor svarar til eit tilsyn med makt til å straffe. Designaren svarar i siste instans til oppdragsgjevaren og til marknaden, og om han lagar noko skadeleg, finst det sjeldan noko fagorgan som kan ta yrket frå han. Det finst kodeksar, men dei manglar tenner. Det finst etisk tenking i mengd — Fry, Escobar, Papanek før dei — men ikkje éi felles, forpliktande etikk som heile faget er samla om og bunde av.

Eg skal ikkje vri meg unna kor mykje dette veg. Eit fag som formar verda — og eg har nett brukt eit heilt kapittel på å hevde at design gjer nett det — og som samstundes manglar bindande handheving, har eit reelt og alvorleg ope rom. Krafta og ansvaret står ikkje i stil. Det er den beste innvendinga motparten har, og ho fortener ikkje eit lettvint svar. Ho fortener at eg seier: ja, dette manglar faget, og fråvêret er ikkje smått.

\section{Kvifor mangelen ikkje er ein dom}

Men her skil vegane lag, og skilnaden ligg i kva ein sluttar av mangelen. \textsc{Grunnlaust} les fråvêret av handheving og samla etikk som prov på at faget er fundamentlaust — at det aldri nådde fram, og at det difor er ein avleggar i vente. Eg les det annleis. Eg les det som veikskapane til eit \emph{ungt} fag, og veikskapar av det slaget eit veksande fag kan rette, ikkje teikn på at det manglar grunn å vekse frå.

Tenk på dei modne profesjonane før dei var modne. Medisinen hadde inga bindande handheving før godt inn på 1800-talet; lenge var legeyrket ope for kven som helst, og dei etiske kodeksane som no har tenner, vart bygde stein for stein over generasjonar, ofte i kjølvatnet av skandalar og strid. Jussen, ingeniørfaget, revisjonen — alle fekk tilsyna og standsorgana sine seint, og fekk dei fordi nokon kjempa dei fram, ikkje fordi faga var fundamentlause til orgena kom og velgrunna dagen etter. Handheving er ikkje fundamentet til eit fag. Det er ei \emph{institusjonell mogning} som kjem oppå fundamentet, når faget er gammalt nok og samfunnet har bestemt at det treng tilsynet. Designfaget i si profesjonelle form er ungt — yngre enn dei fleste samanliknar det med — og at det enno ikkje har vakse dei harde institusjonane, er det ein ventar av eit fag på dette steget, ikkje eit prov på at det står på sand.

Det same gjeld den manglande samla etikken. At faget enno strir om etikken sin — at Fry, Escobar, Manzini og dei andre dreg i ulike retningar — kan lesast som oppløysing, eller det kan lesast som det det ser ut som i alle levande fag: ein open, pågåande samtale om kva faget skuldar verda. Modne fag har ikkje slutta å vere usamde om etikken sin; bioetikken er eit stridsfelt, ikkje eit ferdig kapittel. Skilnaden er at design ikkje endå har samla seg om ein delt minstestandard og gjeve han forpliktande kraft. Det er ein veikskap. Men det er veikskapen til eit fag som framleis byggjer, ikkje til eit fag som har gjeve opp å byggje.

\section{Måten å lese eit ungt fag på}

Korleis veit vi, då, om vi ser eit ungt fag eller ein avleggar? Vi kan ikkje avgjere det ved definisjon, og det var poenget mitt heilt frå starten. Vi må sjå kva faget faktisk gjer. Og det design gjer, er ikkje det ein avleggar gjer. Ein avleggar smuldrar; han misser sjølvtillit, sluttar å produsere kunnskap, blir oppslukt av nabofaga. Designforskinga gjer det motsette. Ho veks. Nye doktorgradsmiljø, nye tidsskrift, ein eigen litteratur som korrigerer seg sjølv — Norman som finskil affordans frå signifier i ny utgåve, Cross som byggjer ut «designerly ways of knowing» til ein heil epistemologi, Nelson og Stolterman som freistar å skrive fram ein eigen «design way» \autocite{nelson2012designway}, Redström som spør korleis designteori i det heile lèt seg lage og kva slags teori eit slikt fag treng \autocite{redstrom2017making}. Dette er rørslene til eit fag som modnar, ikkje til eit som forfell. Eit fag som spør, med Redström, kva si eiga teoriform skal vere, er ikkje ferdig — men det er nett uferdigheita til noko levande, som veit at det har meir å byggje.

Og det byggjer i fleire retningar samstundes. Escobar og dei andre dreg faget mot ein etikk og ein politikk for jorda og for rettferda \autocite{escobar2018pluriverse} — ikkje fordi faget har funne svaret, men fordi det har teke spørsmålet på alvor. Eit ungt fag kjenner ein på at det stiller dei store spørsmåla før det har dei store svara. Det er ikkje skam. Det er ungdom.

\section{Dei to bøkene saman}

Så kvifor skreiv eg begge? Kvifor sette eg meg ned og bygde det sterkaste eg kunne av ei sak eg så vender meg mot i den andre boka? Fordi sanninga om designfaget ikkje ligg i nokon av bøkene åleine. Ho ligg i striden mellom dei.

\textsc{Grunnlaust} er sann om fråværa — om dei manglande institusjonane, den usamla etikken, den vanskelege evna til å seie nei. \textsc{Grunnfest} er sann om grunnlaget — om kunnskapsforma, metoden, den ontologiske krafta, mennesket som feste. Begge bøkene les den same litteraturen; den eine les henne mot håra, den andre med. Og den lesaren som har følgt begge til endes, sit ikkje att med ein vinnar utdelt frå mi hand. Han sit att med noko ærlegare: biletet av eit fag i strid med seg sjølv om sitt eige grunnlag.

Det er, trur eg, det sannaste portrettet eg kunne gje. For det \emph{er} eit fag i slik strid. Designfaget kranglar med seg sjølv om det har eit fundament, om det er vitskap eller kunst eller noko tredje, om det kan seie nei, om det skuldar verda noko og kva. Den krangelen er ikkje teikn på at faget er tomt. Den er teikn på at det er ungt nok til endå å spørje kven det er, og alvorleg nok til ikkje å nøye seg med eit lettvint svar. Eit fag utan grunn kranglar ikkje om grunnen sin; det har ingen å krangle om. Berre eit fag som anar at det \emph{har} ein grunn, men ikkje endå har fått full greie på han, kan stri slik design strir.

Eg har ført forsvaret. Eg meiner det held — at design har eit fundament, sjølv om det ikkje har fysikken sitt, og at fråværa som står att, er veikskapane til eit veksande fag og ikkje dommen over eit dødt eitt. Men eg skreiv motboka òg, og eg står ved henne. Lat dei to bli ståande side om side. Eit fag som toler å bli lese så hardt mot seg sjølv som \textsc{Grunnlaust} les design, og som likevel reiser seg att så sterkt som eg har freista å la det reise seg her, er ikkje ein avleggar i vente. Det er eit ungt fag som veks — i strid, i tvil, og på ein grunn det enno held på å forstå.
